\documentclass[11pt]{article}
\usepackage{amsmath, amssymb}
\usepackage{geometry}
\usepackage{array}
\usepackage{booktabs}
\geometry{margin=0.7in}
\setlength{\parindent}{0pt}
\setlength{\parskip}{0.35em}
\renewcommand{\arraystretch}{1.15}

\newcommand{\cheatsheettitle}[2]{%
\begin{center}
{\LARGE #1}\\[0.4em]
{\large #2}
\end{center}
}


\begin{document}
\cheatsheettitle{Algebra II Polynomial Functions Cheatsheet}{Module 4}

\section*{End Behavior}
Use degree and leading coefficient.
\[
\begin{array}{lll}
\text{even degree} & \text{positive leading coefficient} & \text{both ends up}\\
\text{even degree} & \text{negative leading coefficient} & \text{both ends down}\\
\text{odd degree} & \text{positive leading coefficient} & \text{left down, right up}\\
\text{odd degree} & \text{negative leading coefficient} & \text{left up, right down}
\end{array}
\]

\section*{Zeros and Multiplicity}
If
\[
f(x)=a(x-r)^m(\cdots)
\]
then $r$ is a zero with multiplicity $m$.
\[
m\text{ even}\Rightarrow \text{touches and turns}
\]
\[
m\text{ odd}\Rightarrow \text{crosses}
\]

\section*{Rational Root Theorem}
If
\[
f(x)=a_nx^n+\cdots+a_0
\]
has rational root $\frac pq$ in lowest terms, then
\[
p\mid a_0,\qquad q\mid a_n
\]
Possible rational roots:
\[
\pm\frac{\text{factors of constant term}}{\text{factors of leading coefficient}}
\]

\section*{Remainder and Factor Theorems}
Remainder Theorem:
\[
\text{remainder on division by }x-r\text{ is }f(r)
\]
Factor Theorem:
\[
f(r)=0\iff x-r\text{ is a factor}
\]

\section*{Division}
Use synthetic division only for divisors of the form
\[
x-r
\]
Use long division for other polynomial divisors.

\section*{Complex Roots}
For real-coefficient polynomials:
\[
a+bi\text{ root}\Rightarrow a-bi\text{ root}
\]
Nonreal roots come in conjugate pairs.

\section*{Graph Checklist}
\begin{itemize}
  \item Find real zeros and multiplicities.
  \item Find \(y\)-intercept $f(0)$.
  \item Determine end behavior.
  \item Use turning points and sign checks to refine shape.
\end{itemize}

\end{document}
